% id: 105-198
% book: 개념원리 기하 · 13 정사영
% section: 개념원리 익히기
% figure: none
% answer: ⑴ $6$ ⑵ $8\sqrt{2}$ ⑶ $30^\circ$
% answer_source: 답지
% variant_level: 0 (원본 전사)
% 이 파일은 build.mjs 가 items.json 에서 만든다. 고칠 때는 items.json 을 고친다.
\smash{\raisebox{1.5mm}{\dmkichul{개념원리 기하 105쪽 198번}}}
평면~$\alpha$ 위에 있는 도형의 넓이를 $S$, 이 도형의 평면~$\beta$ 위로의 정사영의 넓이를 $S'$이라 하고, 두 평면~$\alpha$, $\beta$가 이루는 각의 크기를 $\theta$라 할 때, 다음을 구하시오.
\dmsubprob{1} $S=12$, $\theta=60^\circ$일 때, $S'$의 값
\dmsubprob{2} $S'=8$, $\theta=45^\circ$일 때, $S$의 값
\dmsubprob{3} $S=4\sqrt{3}$, $S'=6$일 때, $\theta$의 값
